% !PRINTED 224
\weberpage{224}

%<<CONT
for a substitution of the second order
\weq{15}{a + d = 0,}
and for a substitution of the third order
\weq{16}{a^2 + bc + ad + d^2 = 0.}

Now the substitutions (11), apart from the identical one contained among them, are of the fifth order. The five substitutions (12) are of the second order. Of the substitutions (13), by (15), the five contained in the form $\Theta^r\chi\Theta^{-r}$ (and only these) are of the second order, and consequently five more of the substitutions (14) must also be of the second order. But by (15) this is only possible \gesperrt{if $\beta$ is a power of $\varepsilon$}.

The substitutions of the third order must now distribute themselves over the forms (13) and (14). For these the conditions resulting from (16) are:

that (13) be of the third order
\weq{17}{\omega^2(\varepsilon^{r+s} + \varepsilon^{-(r+s)} - 1) = -\beta,}
and that (14) be of the third order
\weq{18}{\omega^2 = \beta - \varepsilon^{-(r+s)} - \beta^2\varepsilon^{r+s}.}

But if either of the two conditions (17) or (18) is to be satisfiable, we must have:
\[
\beta = 1
\]
For $\omega^2$ is real, and the whole left side of (17) is therefore likewise real; consequently $\beta$ must be real, and, since it is a power of $\varepsilon$, must be equal to 1.

But if (18) is to be satisfied, the left side must not change when one passes to the conjugate imaginary quantities; that is, we must have
\[
\beta - \varepsilon^{-(r+s)} - \beta^2\varepsilon^{r+s} = \beta^{-1} - \varepsilon^{r+s} - \beta^{-2}\varepsilon^{-(r+s)}
\]
or
\[
(\beta - \beta^{-1})(1 - \beta^{-1}\varepsilon^{-(r+s)} - \beta\varepsilon^{r+s}) = 0
\]
and this, because $\beta$ is a power of $\varepsilon$, is possible only if $\beta = \beta^{-1}$, and hence $\beta = 1$.

Accordingly five of the substitutions (14) are of the second order, namely $\Theta^r\chi\psi\Theta^{-r}$, and it follows that (13) is of the third order if
\[
\omega^2(\varepsilon^{r+s} + \varepsilon^{-(r+s)} - 1) = -1,
\]
%CONT>>
